\documentclass[11pt,a4paper]{article}
\usepackage[T1]{fontenc}
\usepackage{lmodern,amsmath,amssymb,amsthm,booktabs,microtype,hyperref}
\usepackage[margin=27mm]{geometry}
\setlength{\emergencystretch}{3em}
\newcommand{\Q}{\mathbb Q}\newcommand{\Z}{\mathbb Z}
\newcommand{\Tr}{\operatorname{Tr}}
\newtheorem{lemma}{Lemma}
\newtheorem{theorem}[lemma]{Theorem}
\newtheorem{proposition}[lemma]{Proposition}
\title{Even-valuation abscissa tests in the general detector}
\author{}\date{17 November 2026}
\begin{document}\maketitle
\begin{abstract}
The first ten field tests in the general detector are reciprocal
abscissas. Their valuations are even outside the retained exceptional
set, so their detector factors do not require squared numerators.
The resulting local implication is stronger, and the original
prepared witnesses retain all real and whole-character guards.
With the full auxiliary charts this gives a complete twenty-five-variable
family, under the unchanged arithmetic premises, of degree at most
$4863007848923136$ and support strictly below $2^{1192}$.
The same change is valid in the normalized two-actual-trace source;
its simultaneous rational realization remains unproved.
\end{abstract}

\section{The revised factors}

Retain the general source of \cite{General,Unit}, its full auxiliary
charts from \cite{FullAux}, and all of their arithmetic premises.
In particular
\[
 K=\Q(\sqrt2),\qquad D=11-6\sqrt2,\qquad
 E:y^2=x^3-Dx,
\]
and the four root points retain their integer indices in $2E(K)$.
The five auxiliary points range over guarded affine points of $E(K)$.
Write $U_0=H$, $Q_0=P$ and keep the other four root and numerator
roles $Q_i,U_i$. The first ten field tests are
\[
 C_1,\ldots,C_{10}
 =\bigl(f(U_i)\bigr)_{i=0}^4,
   \bigl(\overline{f(U_i-Q_i)}\bigr)_{i=0}^4,
 \qquad f(Q)=1/x(Q).
\]
They have order one. The six remaining field tests, their orders
$a_i\in\{1,3\}$, the six rational scalar tests $S_k$, their orders
$\kappa_k\in\{1,3\}$, the contact functions $R_j$, normalizers
$\delta_j$ and scalar $\rho$ are unchanged.

For the compiled family put $q_j=\rho R_j$ and retain the original
scalar normalizations
\[
 G_{j,k}=R_j^{\kappa_k}\frac{R_jS_k}{1+R_jS_k}.
\]
Replace the detector product by
\begin{equation}\label{eq:product}
 W_j=
 \prod_{i=1}^{10}\left(1+\frac{C_i}{\delta_j^2q_j}\right)
 \prod_{i=11}^{16}
       \left(1+\frac{C_i^2}{\delta_j^2q_j^{2a_i-1}}\right)
 \prod_{k=1}^{6}
       \left(1+\frac{G_{j,k}^2}{q_j^{2\kappa_k-1}}\right).
\end{equation}
Keep $V_j=\Tr_{K/\Q}(W_j^{-1})$, the two native pole inputs
$b_j=4q_j/V_j^2$, and every source guard. The only native circuit
change is the power of each of the first ten field-test numerators.

\section{Even valuations and local extraction}

\begin{lemma}\label{lem:even}
Let $v$ be an integer-valued nonarchimedean valuation of a
characteristic-zero field, with $v(D)=0$. Every finite point of
$y^2=x^3-Dx$ with nonzero abscissa has even $v(x)$ and hence
even $v(1/x)$.
\end{lemma}
\begin{proof}
Put $m=v(x)$. If $m>0$, the two terms on the right have distinct
valuations $3m$ and $m$, so $2v(y)=m$. If $m<0$, the smaller
valuation is $3m$, so $2v(y)=3m$. Either equality makes $m$
even. The remaining case is $m=0$.
\end{proof}

The norm of $D$ is $49$. At every rational prime outside $2,7$,
the extension $K/\Q$ is unramified and $D$ is a local unit.
Thus each of the first ten tests has even valuation at every place
above every consumed prime. The retained exceptional set already
contains $2$ and $7$. No restriction to the subgroup $2E(K)$ is
needed for this local assertion.

\begin{proposition}\label{prop:local}
Suppose the unchanged native pole test excludes a negative odd
valuation of $b_j$ at a nonexceptional prime. If
$\alpha=v_p(q_j)$ is odd, it is positive. At some place above $p$
every factor of \eqref{eq:product} has valuation greater than
$-\alpha/2$. If $\delta_j$ is a unit there, then
\[
 v(C_i)>\alpha/2\quad(1\le i\le10),\qquad
 v(C_i)>(a_i-3/4)\alpha\quad(11\le i\le16),
\]
and the scalar inequalities of the original source are unchanged.
\end{proposition}
\begin{proof}
For each revised factor the valuation of its second summand is
$v(C_i)-2v(\delta_j)-\alpha$, an odd integer by
Lemma~\ref{lem:even}. The second summands of the other factors
likewise have odd valuation when nonzero. Thus every factor has
nonpositive valuation, with no cancellation at valuation zero.
A zero test contributes the factor one.

Let $m$ be the larger of the two valuations of $W_j$. Then
$m\le0$ and integrality of local traces gives $v_p(V_j)\ge-m$.
If $\alpha<0$, the odd number $v_p(b_j)$ is at most
$\alpha+2m<0$, a contradiction. Otherwise the pole implication
gives
\[
             1\le v_p(b_j)\le\alpha+2m,
             \qquad m>-\alpha/2.
\]
At a place attaining $m$, each factor has valuation at least the
sum $m$. Applying this to a revised summand at a unit normalizer
gives $v(C_i)-\alpha>-\alpha/2$. The unchanged squared factors
give the other inequalities exactly as before.
\end{proof}

In particular the new first-ten inequality implies the old one
$v(C_i)>\alpha/4$. All scalar-odd and mixed-contact alternatives,
formal logarithmic comparisons, integer residual divisibilities,
and subsequent height arguments therefore survive. The full
auxiliary charts satisfy the same formal-index conclusions proved
in \cite{FullAux}.

\section{Completeness and guards}

Use the original ordered construction of prepared witnesses, with
all chosen points in $2E(K)$. At a selected fresh contact prime,
$v_p(q_j)=1$, $\delta_j$ is a unit and every successful-slot
reciprocal abscissa has valuation at least two. Each revised
summand consequently has valuation at least one, so its factor
is congruent to one. At the opposite slot its valuation is odd;
the factor is either congruent to one or has negative valuation.
The reciprocal product there is therefore zero or one modulo the
selected place. The scalar and other comparison factors are
unchanged. The trace remains a unit, as required for the positive
odd pole-input valuation.

At the fixed preparation places, the scalar preparation is
unchanged. After fixing the compact point neighborhoods, the first
ten raw tests are bounded there, and their revised summands tend to
zero as $\delta_j^{-1}$ tends to zero and $R_j$ tends to one.
Thus the prescribed arbitrary finite precisions and real bounds
can still be imposed before selecting the slope and fresh primes.
The same denominator-square and native pole-input preparation
arguments apply.

For the whole-character real conditions, retain the old complete
subfamily inside the enlarged point charts. At each real embedding,
its points and guarded differences lie in $2E(\mathbb R)$, the
identity component of $E(\mathbb R)$. Since $D>0$ at both
embeddings, a finite point in this component has
$x\ge\sqrt D>0$. Hence each revised $C_i$ is positive at every
prepared character. With $q_j>0$ and $\delta_j\ne0$, its factor
is at least one. All product and reciprocal-trace norm guards are
therefore preserved. The old and enlarged auxiliary point algebras
agree at these witnesses by \cite{FullAux}; no new denominator or
unrestricted coordinate has been added.

\begin{theorem}\label{thm:complete}
Under exactly the retained arithmetic premises of the general
source, there is a computable family of integer polynomials in
twenty-five total rational coordinates, including the input, whose
input projection is $\Z$ for every sufficiently large fixed odd
$L$. Its degree is at most $4863007848923136<2^{53}$, its support
is strictly below $2^{1192}$, and its integer content is one.
The sufficient defining exponent remains ineffective.
\end{theorem}
\begin{proof}
Proposition~\ref{prop:local} supplies the original local
implications and hence soundness for every sufficiently large fixed
odd $L$. The preceding prepared construction gives completeness for
every integer input and every fixed odd $L\ge3$. The two pole blocks
and the finite-algebra separator are unchanged. Their count is
$1+18+1+4+1=25$, including the input.

The exact circuit package below retains the point-basis and terminal
root bounds of \cite{FullAux}. It computes
\[
 \deg G_*\le4637555553,\qquad \deg\Omega\le170132,
\]
so the rank-$1048576$ encoder has degree at most
\[
 1048576(4637555553+170132+1)=4863007848923136.
\]
The encoder-coordinate degree is at most $1048576$. Counting
monomials under these two bounds gives support strictly below
$2^{1192}$. Setting the encoder coordinate to zero gives one,
so the integer content is one.
\end{proof}

\section{The normalized actual-trace source}

The same first-ten change is valid in the reweighted general source
of \cite{Reweighted}. There use $w_j=q_j/(1+q_j)^2$ in place
of $q_j$ throughout \eqref{eq:product}, and keep that source's
scalar numerators $R_j^{\kappa_k+1}\Phi_{R_j}(S_k)$.
The parity argument is identical and the new extraction is again
stronger. Selected contacts still have $v_p(w_j)=1$; at fixed
preparation places $w_j$ tends to the same nonzero limiting value.
Thus this revision preserves the complete twenty-five-variable
normalized source and the sound twenty-one-variable actual-trace
source with the same guards and local preparation. The numerical
degree bound in Theorem~\ref{thm:complete} is for the compiled
unnormalized family \eqref{eq:product} only.

No simultaneous rational square realization follows. At fixed
guarded point values, each of the first detector's ten contacts
is still a simple zero of $w_0$. In each slot its nonzero
reciprocal tests give poles of $W_0$, while the scalar factors tend
to one. Hence
\[
 V_0=O(w_0),\qquad
 V_0^2+w_0(2-V_0)=2w_0+O(w_0^2).
\]
The ten simple branch points and genus-at-least-four obstruction
of \cite{Obstruction} therefore remain. In particular this
simplification does not create a rational parametrization of the
fixed-point square cover or settle the complete twenty-one-variable
target.

At a fixed input and fixed $L$, the scalar $\rho$ is fixed by the
source. The other scale is $1+\rho T^2$, with
$T=\Tr\sum_i x(Q_i)$ determined by the four integer-index root
points. Neither scalar is an independent free rational parameter.
The odd valuation of $\rho$ at $76249$ excludes the value
$1+\rho T^2=2$. An arbitrary auxiliary square rescaling of $\rho$
would also change its height estimate and is not an operation
justified by the retained proof.

\section{Exact specification and checks}

The package
\begin{quote}\small
\nolinkurl{computation/few_variables/even_valuation25_v75}
\end{quote}
specifies the integer polynomial family by its native circuit and
finite-algebra determinant. It has $6039$ integer arithmetic nodes,
$101$ recorded guards and $128$ retained norm factors. The stored
$L=3$ member is an algebraic fixture; its coefficient $\ell^1$-norm
bound is $2^{34532878499971072}$. This value of $L$ is not
asserted to exceed the sufficient defining threshold.

The independent checker verifies the point-algebra identities,
guard and subalgebra supports, weighted degrees and coefficient
bounds, and rejects deliberate degree and guard corruption. It
independently evaluates the revised native detector at $72$ finite
characters of actual point charts, including half-integral auxiliary
indices. The direct evaluator uses unsquared first-ten factors
without reading the circuit's product expression. These finite
checks certify the algebraic implementation and bounds; they do
not prove the retained arithmetic premises or global actual-square
realization.

\begin{thebibliography}{9}\raggedright
\bibitem{General}\emph{Twenty-five variables from combined scalar and
comparison conditions},
\nolinkurl{papers/few_variables25/few_variables25_merged_20261110_v19.tex}.
\bibitem{Unit}\emph{A single norm determinant for the 25- and
13-variable constructions},
\nolinkurl{papers/finite_algebra/unit_norm_20261114_v28.tex}.
\bibitem{FullAux}\emph{Full auxiliary point charts for the general
twenty-five-variable source},
\nolinkurl{papers/few_variables25/full_auxiliary25_20261117_v72.tex}.
\bibitem{Reweighted}\emph{Quadratic reweightings of native elliptic
detectors},
\nolinkurl{papers/pole_parameters/reweighted_elliptic_detectors_20261116_v53.tex}.
\bibitem{Obstruction}\emph{Curves and abelian families in an
actual-trace square locus},
\nolinkurl{negative_results/papers/pole_parameters/curves_and_abelian_families_20261117_v69.tex}.
\end{thebibliography}
\end{document}
